% =====================================================================
% Chủ đề 4 — Số học và Tư duy số — phần a: lý thuyết, ví dụ, bài 1–15
% =====================================================================
\begin{knowbox}
\textbf{Ý tưởng lớn.} ``Tư duy số'' (cảm nhận về số) là khả năng \emph{hiểu} con số chứ không chỉ \emph{tính} con số:
biết số nào lớn hơn, một kết quả có hợp lý không, nên làm tròn thế nào để ước lượng nhanh, và nhìn ra cấu trúc ẩn
bên trong một số (các chữ số, các ước, các bội). Trong chủ đề này, em sẽ ôn lại và nối kết các mảng kiến thức:
\begin{itemize}[leftmargin=1.4em,itemsep=1pt,topsep=2pt]
  \item \textbf{số thập phân và số nguyên âm}: so sánh, sắp xếp trên trục số;
  \item \textbf{thứ tự thực hiện phép tính} và \textbf{ước lượng} kết quả;
  \item \textbf{phân số, tỉ số, tỉ số phần trăm}: ``một phần của một số'', phân số bằng nhau, giảm giá;
  \item \textbf{tốc độ} -- quãng đường -- thời gian và đổi đơn vị (km/h, m/s);
  \item \textbf{chia hết, ước chung, chữ số tận cùng, cấu tạo số} (giải các phép tính đố chữ);
  \item \textbf{trung bình cộng, trung vị} và các \textbf{dãy số có quy luật} (Fibonacci, cấp số cộng).
\end{itemize}
Bí quyết chung: \emph{trước khi tính, hãy ước lượng; sau khi tính, hãy thử lại}. Một học sinh có tư duy số tốt
thường loại được ngay vài phương án vô lý trong câu trắc nghiệm mà chưa cần tính chính xác.
\end{knowbox}
\vspace{0.3cm}

% ---------------------------------------------------------------------
\shead{Số thập phân: đọc, so sánh, sắp xếp, cộng trừ}
\begin{knowbox}
\textbf{1. Giá trị theo hàng.} Trong số thập phân, mỗi chữ số đứng ở một \emph{hàng}. Sau dấu phẩy lần lượt là
hàng phần mười, hàng phần trăm, hàng phần nghìn, \ldots
\begin{center}
\renewcommand{\arraystretch}{1.25}
\begin{tabular}{|c|c|c|c|c|c|}
\hline
\rowcolor{tmBlue!10} Chục & Đơn vị & , & Phần mười & Phần trăm & Phần nghìn\\ \hline
 & $2$ & , & $2$ & $3$ & \\ \hline
 & $2$ & , & $3$ & $0$ & \\ \hline
\end{tabular}
\end{center}
Ví dụ: $2{,}23 = 2 + \dfrac{2}{10} + \dfrac{3}{100}$, còn $2{,}3 = 2 + \dfrac{3}{10} = 2{,}30$.

\textbf{2. Thêm hoặc bớt chữ số $0$ ở tận cùng phần thập phân không làm thay đổi giá trị:}
$2{,}3 = 2{,}30 = 2{,}300$. Lý do: $\dfrac{3}{10} = \dfrac{30}{100} = \dfrac{300}{1000}$ (nhân cả tử và mẫu với $10$).

\textbf{3. Quy tắc so sánh hai số thập phân dương.}
\begin{itemize}[leftmargin=1.4em,itemsep=1pt,topsep=2pt]
  \item So sánh \emph{phần nguyên} trước: số nào có phần nguyên lớn hơn thì lớn hơn.
  \item Nếu phần nguyên bằng nhau, \emph{viết thêm chữ số $0$} để hai phần thập phân có cùng số chữ số, rồi so sánh
  từng hàng từ trái sang phải: hàng phần mười, rồi hàng phần trăm, \ldots
\end{itemize}
Ví dụ: so sánh $0{,}304$ và $0{,}34$. Viết $0{,}34 = 0{,}340$. Hàng phần mười: $3 = 3$; hàng phần trăm: $0 < 4$.
Vậy $0{,}304 < 0{,}340 = 0{,}34$.

\textbf{4. Cộng, trừ số thập phân:} đặt tính sao cho \emph{các dấu phẩy thẳng cột} (hàng thẳng hàng), rồi cộng/trừ
như số tự nhiên.

\textbf{5. Làm tròn số thập phân.} Muốn làm tròn đến hàng phần trăm (hai chữ số sau dấu phẩy), ta nhìn chữ số ngay
bên phải hàng đó (hàng phần nghìn): nếu chữ số này nhỏ hơn $5$ thì giữ nguyên, nếu lớn hơn hoặc bằng $5$ thì cộng thêm
$1$ vào hàng phần trăm; sau đó bỏ các chữ số phía sau. Ví dụ: $4{,}6875 \approx 4{,}69$ (vì chữ số hàng phần nghìn là
$7 \ge 5$), còn $3{,}1416 \approx 3{,}14$ (vì $1 < 5$). Kí hiệu ``$\approx$'' đọc là ``xấp xỉ''.
\end{knowbox}
\begin{proofbox}
Vì sao không được so sánh ``theo độ dài'' (số nào nhiều chữ số hơn thì lớn hơn)? Hãy viết theo phần nghìn:
$0{,}304 = \dfrac{304}{1000}$ và $0{,}34 = \dfrac{340}{1000}$. Hai phân số cùng mẫu, tử $304 < 340$ nên
$0{,}304 < 0{,}34$, dù $0{,}304$ ``trông dài hơn''. Việc thêm chữ số $0$ chính là đưa về cùng mẫu số.
\end{proofbox}
\begin{warnbox}
Nhiều bạn nghĩ $2{,}23 > 2{,}3$ vì ``$23 > 3$''. Sai! Phải so sánh từng hàng: $2{,}23$ có $2$ phần mười,
còn $2{,}3$ có $3$ phần mười, nên $2{,}23 < 2{,}3$.
\end{warnbox}
\begin{tipbox}
Muốn tìm số nhỏ nhất (lớn nhất) trong một tập hợp số thập phân, hãy viết tất cả các số thành \emph{một cột}, dấu phẩy
thẳng nhau, thêm chữ số $0$ cho đủ dài. Khi đó em chỉ việc so sánh như các số tự nhiên.
\end{tipbox}

% ---------------------------------------------------------------------
\shead{Số nguyên âm và trục số}
\begin{knowbox}
\textbf{1. Trục số.} Các số nguyên được xếp trên một đường thẳng: số $0$ ở giữa, các số dương ở bên phải, các số âm
ở bên trái.
\numlinex{8}{0/$-4$,1/$-3$,2/$-2$,3/$-1$,4/$0$,5/$1$,6/$2$,7/$3$,8/$4$}{}{4}
\textbf{2. Quy tắc so sánh.} Trên trục số (nằm ngang), \emph{số bên trái nhỏ hơn số bên phải}. Do đó:
\begin{itemize}[leftmargin=1.4em,itemsep=1pt,topsep=2pt]
  \item mọi số âm đều nhỏ hơn $0$, và $0$ nhỏ hơn mọi số dương;
  \item mọi số âm nhỏ hơn mọi số dương;
  \item với hai số âm, số nào \emph{xa số $0$ hơn} thì \emph{nhỏ hơn}: $-8 < -2$ vì $8 > 2$.
\end{itemize}
\textbf{3. Hình ảnh đời thường.} Nhiệt độ $-8^\circ$C lạnh hơn $-2^\circ$C; nợ $8$ đô la thì ``nghèo hơn'' nợ
$2$ đô la (số tiền là $-8$ so với $-2$).

\textbf{4. Mặt phẳng tọa độ.} Một điểm được xác định bởi cặp số $(x;y)$: $x$ là hoành độ (đi sang phải/trái),
$y$ là tung độ (đi lên/xuống). Di chuyển ngang chỉ làm thay đổi $x$, di chuyển dọc chỉ làm thay đổi $y$.
\end{knowbox}
\begin{proofbox}
Vì sao $-8 < -2$? Ta có $-2 = -8 + 6$: muốn đi từ $-8$ đến $-2$ trên trục số ta phải đi \emph{sang phải} $6$ đơn vị.
Đi sang phải nghĩa là tăng lên, nên $-2$ lớn hơn $-8$.
\end{proofbox}
\begin{warnbox}
Đừng so sánh số âm ``theo cảm giác của số dương''. $-8$ \emph{nhỏ hơn} $-2$, dù $8$ lớn hơn $2$. Hãy luôn vẽ
nhanh một trục số khi phân vân.
\end{warnbox}
\begin{tipbox}
\textbf{Tính chẵn -- lẻ} là một công cụ ``tư duy số'' rất mạnh: chẵn $+$ chẵn $=$ chẵn, lẻ $+$ lẻ $=$ chẵn, chẵn $+$ lẻ $=$ lẻ;
tích của các số lẻ luôn là số lẻ. Chẳng hạn, tổng của $5$ số lẻ là số lẻ, tổng của $4$ số lẻ là số chẵn. Khi một điểm
di chuyển trên lưới tọa độ, xét tính chẵn -- lẻ của tọa độ thường cho biết ngay một vị trí có thể đến được hay không.
\end{tipbox}

% ---------------------------------------------------------------------
\shead{Thứ tự thực hiện phép tính và ước lượng}
\begin{knowbox}
\textbf{1. Thứ tự thực hiện phép tính.}
\begin{enumerate}[leftmargin=1.6em,itemsep=1pt,topsep=2pt]
  \item Trong \emph{ngoặc} trước: $(\;)$, rồi $[\;]$, rồi $\{\;\}$.
  \item Tiếp theo là \emph{lũy thừa}.
  \item Rồi đến \emph{nhân và chia}, làm lần lượt từ trái sang phải.
  \item Cuối cùng là \emph{cộng và trừ}, làm lần lượt từ trái sang phải.
\end{enumerate}
Ví dụ: $6\times 2 - 3 = 12 - 3 = 9$ (nhân trước, trừ sau), còn $6\times(2-3) = 6\times(-1) = -6$.

\textbf{2. Làm tròn để ước lượng.} Để biết nhanh một kết quả ``khoảng bao nhiêu'', ta làm tròn các thừa số đến số
tròn chục, tròn trăm gần nhất rồi tính nhẩm. Ví dụ: $99 \times 9 \approx 100 \times 9 = 900$.

\textbf{3. Mẹo tính nhanh:} $99 \times 9 = (100-1)\times 9 = 900 - 9 = 891$ (tính chất phân phối của phép nhân
đối với phép trừ).
\end{knowbox}
\begin{proofbox}
Vì sao $(100-1)\times 9 = 100\times 9 - 1\times 9$? Hãy tưởng tượng $9$ hàng, mỗi hàng $100$ chiếc kẹo,
nhưng mỗi hàng bị lấy đi $1$ chiếc. Số kẹo còn lại là $9$ hàng $\times$ $99$ chiếc, cũng bằng
($9\times 100$ chiếc) trừ đi ($9$ chiếc bị lấy). Đó là \emph{tính chất phân phối}: $a\times(b-c) = a\times b - a\times c$.
\end{proofbox}
\begin{warnbox}
Tính ``từ trái sang phải bất kể phép tính gì'' là sai: $2 + 3\times 4 \ne 5\times 4$. Đúng là
$2 + 3\times 4 = 2 + 12 = 14$.
\end{warnbox}
\begin{tipbox}
Trong câu trắc nghiệm hỏi ``số nào gần nhất với\ldots'', các phương án thường cách nhau gấp $10$ lần
($10$; $100$; $1000$; \ldots). Khi đó chỉ cần ước lượng thô là đủ chọn đúng.
\end{tipbox}

% ---------------------------------------------------------------------
\shead{Phân số: một phần của một số, phân số bằng nhau}
\begin{knowbox}
\textbf{1. Phân số là một phần của cái toàn thể.} $\dfrac{a}{b}$ nghĩa là chia toàn thể thành $b$ phần bằng nhau
và lấy $a$ phần.

\textbf{2. Tìm giá trị phân số của một số:} muốn tìm $\dfrac{m}{n}$ của $a$, ta tính $a \times \dfrac{m}{n}$
(chia $a$ thành $n$ phần rồi lấy $m$ phần). Ví dụ: $\dfrac{1}{3}$ của $396$ là $396 : 3 = 132$.

\textbf{3. Phần còn lại.} Nếu đã làm được $\dfrac{m}{n}$ công việc thì phần còn lại là $1 - \dfrac{m}{n} = \dfrac{n-m}{n}$.

\textbf{4. Phân số bằng nhau (tính chất cơ bản).} Nhân (hoặc chia) cả tử và mẫu với cùng một số khác $0$ thì được
phân số bằng nó:
\[ \frac{3}{7} = \frac{3\times 9}{7\times 9} = \frac{27}{63}, \qquad \frac{12}{36} = \frac{12:12}{36:12} = \frac{1}{3}. \]
\textbf{5. Rút gọn phân số:} chia cả tử và mẫu cho ước chung lớn nhất của chúng để được phân số tối giản.

\textbf{6. Hai số biết tỉ số và tổng (hoặc hiệu).} Nếu $\dfrac{x}{y} = \dfrac{5}{8}$ thì có một số $k$ sao cho
$x = 5k$ và $y = 8k$ (vì $\dfrac{5k}{8k} = \dfrac58$). Biết tổng $x + y$ (hoặc hiệu $y - x$), ta lập phương trình
ẩn $k$ rồi giải. Ví dụ: $x + y = 91$ cho $5k + 8k = 91$, tức $13k = 91$, $k = 7$; suy ra $x = 35$, $y = 56$.

\textbf{7. Phân số của một phân số.} ``Một nửa của một phần tư'' là $\dfrac12\times\dfrac14 = \dfrac18$: nhân tử với tử,
mẫu với mẫu. Điều này hay gặp khi một hình được chia đôi nhiều lần liên tiếp: mỗi lần chia đôi, phần diện tích nhân thêm
với $\dfrac12$.

\textbf{8. So sánh phân số.} Quy đồng mẫu số rồi so sánh tử số, hoặc đổi ra số thập phân. Ví dụ:
$\dfrac23 = \dfrac{8}{12} > \dfrac{7}{12}$; \ $\dfrac35 = 0{,}6 < \dfrac34 = 0{,}75$.
\end{knowbox}
\begin{proofbox}
Vì sao nhân cả tử và mẫu với cùng một số thì phân số không đổi? $\dfrac{3}{7}$ là ``chia cái bánh thành $7$ phần,
lấy $3$ phần''. Nếu cắt \emph{mỗi} phần thành $9$ miếng nhỏ hơn, cái bánh có $7\times 9 = 63$ miếng, còn phần ta lấy
có $3\times 9 = 27$ miếng. Lượng bánh không đổi, nên $\dfrac{3}{7} = \dfrac{27}{63}$.
\end{proofbox}
\begin{warnbox}
Đọc kỹ câu hỏi: đề hỏi phần \emph{đã làm} hay phần \emph{còn lại}? Nhiều bạn tính đúng $\dfrac13$ của $396$ là $132$
rồi khoanh luôn, trong khi đề hỏi số trang \emph{còn lại} là $396 - 132 = 264$.
\end{warnbox}
\begin{parentbox}
Ở lớp 7, hãy khuyến khích con \emph{gọi ẩn}: đặt chữ $x$ cho đại lượng chưa biết, viết các điều kiện của đề thành
đẳng thức, giải rồi thử lại. Chẳng hạn ``đã đọc $\dfrac13$ cuốn sách $396$ trang, còn $x$ trang'' viết thành
$132 + x = 396$. Cách này dùng được cho mọi bài, kể cả những bài số liệu ``xấu'' mà hình vẽ không giúp được.
\end{parentbox}

% ---------------------------------------------------------------------
\shead{Tỉ số phần trăm, đơn giá và giảm giá}
\begin{knowbox}
\textbf{1. Phần trăm nghĩa là ``trên một trăm''.} $80\% = \dfrac{80}{100} = 0{,}8$.

\textbf{2. Tìm $p\%$ của một số $a$:} tính $a \times \dfrac{p}{100}$. Ví dụ: $80\%$ của $50$ là
$50 \times \dfrac{80}{100} = 40$.

\textbf{3. Tìm tỉ số phần trăm của $a$ so với $b$:} tính $\dfrac{a}{b}\times 100\%$. Ví dụ: đúng $32$ câu trên $40$ câu
là $\dfrac{32}{40}\times 100\% = 80\%$.

\textbf{4. Đơn giá.} Đơn giá là giá của \emph{một} đơn vị (một ki-lô-gam, một quyển, \ldots):
\[ \text{Thành tiền} = \text{đơn giá} \times \text{số lượng}, \qquad \text{đơn giá} = \text{thành tiền} : \text{số lượng}. \]
\textbf{5. Giảm giá $p\%$.} Giá mới $= $ giá cũ $\times \left(1 - \dfrac{p}{100}\right)$. Giảm $10\%$ tức là còn
$90\%$: nhân với $0{,}9$.

\textbf{6. Giảm giá nhiều lần liên tiếp.} Mỗi lần giảm $10\%$ \emph{của giá hiện tại} thì ta nhân thêm một lần với
$0{,}9$. Ví dụ, giá $100$ đô la sau hai lần giảm: $100 \to 90 \to 81$ (đô la), \emph{không phải} $80$.
\end{knowbox}
\begin{proofbox}
Vì sao hai lần giảm $10\%$ không bằng một lần giảm $20\%$? Lần thứ hai giảm $10\%$ của $90$ (là $9$), chứ không phải
$10\%$ của $100$ (là $10$). Vì giá đã nhỏ đi nên số tiền được giảm lần sau cũng nhỏ đi: tổng cộng chỉ giảm
$10 + 9 = 19$ đô la, tức $19\%$.
\end{proofbox}
\begin{tipbox}
Tìm $10\%$ thật nhanh: chỉ cần chia cho $10$ (dời dấu phẩy sang trái một chữ số). Từ $10\%$ suy ra
$5\%$ (chia đôi), $20\%$ (gấp đôi), $80\% = 100\% - 20\%$, \ldots
\end{tipbox}

% ---------------------------------------------------------------------
\shead{Tốc độ, quãng đường, thời gian và đổi đơn vị}
\begin{knowbox}
\textbf{1. Ba công thức gắn với nhau.} Với chuyển động đều:
\[ s = v \times t, \qquad v = s : t, \qquad t = s : v, \]
trong đó $s$ là quãng đường, $v$ là vận tốc (tốc độ), $t$ là thời gian.

\textbf{2. Ai nhanh hơn?} Người nhanh hơn là người đi được quãng đường \emph{dài hơn} trong \emph{cùng một khoảng thời
gian}. Cách chắc chắn nhất: tính vận tốc của từng người theo cùng một đơn vị (ví dụ m/s) rồi so sánh.

\textbf{3. Chạy nhiều vòng.} Chạy $n$ vòng quanh đường dài $L$ mét thì quãng đường là $n \times L$ mét.

\textbf{4. Đổi km/h sang m/s.} $1$ km $= 1000$ m, $1$ giờ $= 3600$ giây, nên
\[ 1 \text{ km/h} = \frac{1000 \text{ m}}{3600 \text{ s}} = \frac{5}{18} \text{ m/s}, \qquad
   18 \text{ km/h} = 18\times\frac{5}{18} = 5 \text{ m/s}, \qquad 1 \text{ m/s} = 3{,}6 \text{ km/h}. \]
\end{knowbox}
\begin{proofbox}
Vì sao $v = s : t$ và $t = s : v$? Đó chỉ là cách tìm một thừa số chưa biết trong đẳng thức $s = v\times t$:
chia cả hai vế cho $t$ (khác $0$) ta được $v = \dfrac{s}{t}$; chia cả hai vế cho $v$ (khác $0$) ta được $t = \dfrac{s}{v}$.
\end{proofbox}
\begin{warnbox}
Hai đại lượng phải \emph{cùng hệ đơn vị}: nếu quãng đường tính bằng mét và thời gian bằng giây thì vận tốc ra m/s;
đừng chia mét cho giờ. Trước khi so sánh hai vận tốc, hãy đưa về cùng đơn vị.
\end{warnbox}

% ---------------------------------------------------------------------
\shead{Chia hết, ước chung, chữ số tận cùng và cấu tạo số}
\begin{knowbox}
\textbf{1. Bội và ước.} $a$ là \emph{bội} của $b$ (hay $b$ là \emph{ước} của $a$) nếu $a = b\times k$ với $k$ là số
tự nhiên. Ví dụ: $45 = 9\times 5$ nên $45$ là bội của $9$.

\textbf{2. Dấu hiệu chia hết.}
\begin{itemize}[leftmargin=1.4em,itemsep=1pt,topsep=2pt]
  \item cho $2$: chữ số tận cùng là $0, 2, 4, 6, 8$; \quad cho $5$: tận cùng là $0$ hoặc $5$; \quad cho $10$: tận cùng là $0$;
  \item cho $3$: tổng các chữ số chia hết cho $3$; \quad cho $9$: \emph{tổng các chữ số chia hết cho $9$}.
\end{itemize}
\textbf{3. Ước chung.} Số chia hết cả $a$ lẫn $b$ gọi là ước chung của $a$ và $b$. Mọi ước chung đều là ước của
\emph{ước chung lớn nhất} ƯCLN$(a,b)$. Ví dụ: ƯCLN$(40,72)=8$, nên các ước chung của $40$ và $72$ là
$1, 2, 4, 8$.

\textbf{4. Chữ số tận cùng của một tích} chỉ phụ thuộc vào chữ số tận cùng của các thừa số:
$\ldots 3 \times \ldots 7$ tận cùng bằng chữ số tận cùng của $3\times 7 = 21$, tức là $1$.

\textbf{5. Cấu tạo số.} Một số viết bằng các chữ số được phân tích theo hàng:
\[ \nd{abc} = 100a + 10b + c, \qquad \nd{PQQ} = 100P + 10Q + Q = 100P + 11Q. \]
(Dấu gạch ngang trên đầu cho biết đây là số có các chữ số $a, b, c$, không phải tích $a\times b\times c$.)
\end{knowbox}
\begin{proofbox}
Vì sao dấu hiệu chia hết cho $9$ đúng? Xét số $\nd{abc} = 100a + 10b + c$. Ta tách:
\[ 100a + 10b + c = (99a + 9b) + (a + b + c) = 9\times(11a + b) + (a+b+c). \]
Phần $9\times(11a+b)$ luôn chia hết cho $9$. Vậy $\nd{abc}$ chia hết cho $9$ khi và chỉ khi phần còn lại
$a+b+c$ (tổng các chữ số) chia hết cho $9$. Với số có nhiều chữ số hơn, lập luận y hệt vì
$10 = 9+1$, $100 = 99+1$, $1000 = 999+1$, \ldots
\end{proofbox}
\begin{tipbox}
Gặp bài ``đố chữ'' (mỗi chữ cái là một chữ số) như $\nd{PQQ}+\nd{PPQ}+\nd{QQQ}=876$, hãy làm \emph{từ hàng đơn vị
sang trái}, nhớ để ý các số nhớ. Hàng đơn vị thường có ít khả năng nhất, nên tìm được chữ số đầu tiên rất nhanh.
Cũng có thể dùng cấu tạo số để đưa về một phương trình. Khi phải đếm hoặc tìm hết các khả năng (các bộ chữ số,
các kết quả khi tung đồng xu, \ldots), hãy \emph{liệt kê có hệ thống} (theo thứ tự tăng dần, hoặc bằng bảng/sơ đồ
cây) để không sót, không trùng.
\end{tipbox}
\begin{knowbox}
\textbf{Liệt kê có hệ thống và khả năng xảy ra.}
\begin{itemize}[leftmargin=1.4em,itemsep=1pt,topsep=2pt]
  \item Muốn đếm các số có ba chữ số mà \emph{tích} các chữ số bằng một số cho trước, ta tìm các bộ ba chữ số (viết theo
  thứ tự tăng dần để không trùng), rồi đếm số cách sắp xếp mỗi bộ: ba chữ số khác nhau có $3\times 2\times 1 = 6$ cách xếp;
  hai chữ số giống nhau có $3$ cách xếp. Ví dụ: tích bằng $12$ với bộ $\{1;3;4\}$ cho $6$ số $134, 143, 314, 341, 413, 431$;
  bộ $\{2;2;3\}$ cho $3$ số $223, 232, 322$.
  \item Khi mọi kết quả đều \emph{có khả năng như nhau}, xác suất của một sự kiện bằng
  \[ \frac{\text{số kết quả thuận lợi}}{\text{tổng số kết quả có thể}}. \]
  Ví dụ: tung $2$ đồng xu có $4$ kết quả đồng khả năng: SS, SN, NS, NN (S: sấp, N: ngửa). Xác suất cả hai đồng xu cùng
  một mặt là $\dfrac{2}{4} = \dfrac12$. Dùng sơ đồ cây để liệt kê:
\end{itemize}
\begin{figC}
  \node[draw=Navy,circle,inner sep=1.5pt] (r) at (0,0) {\scriptsize bắt đầu};
  \node[draw=tmBlue,fill=tmBlue!10,circle,inner sep=2pt] (a) at (2,0.8) {S};
  \node[draw=tmOrange,fill=tmOrange!10,circle,inner sep=2pt] (b) at (2,-0.8) {N};
  \foreach \n/\y/\l/\c in {aa/1.2/S/tmBlue,ab/0.4/N/tmOrange}{\node[draw=\c,fill=\c!10,circle,inner sep=2pt] (\n) at (4,\y) {\l};}
  \foreach \n/\y/\l/\c in {ba/-0.4/S/tmBlue,bb/-1.2/N/tmOrange}{\node[draw=\c,fill=\c!10,circle,inner sep=2pt] (\n) at (4,\y) {\l};}
  \draw[Slate] (r)--(a) (r)--(b) (a)--(aa) (a)--(ab) (b)--(ba) (b)--(bb);
  \node[right,font=\scriptsize] at (4.3,1.2) {SS \ (cùng mặt)};
  \node[right,font=\scriptsize] at (4.3,0.4) {SN};
  \node[right,font=\scriptsize] at (4.3,-0.4) {NS};
  \node[right,font=\scriptsize] at (4.3,-1.2) {NN \ (cùng mặt)};
\end{figC}
\end{knowbox}
\begin{warnbox}
``Tổng các chữ số chia hết cho $9$'' là dấu hiệu của $9$, \emph{không} dùng cho $4$, $6$ hay $8$. Và một số chia hết
cho $3$ chưa chắc chia hết cho $9$: $12$ chia hết cho $3$ nhưng không chia hết cho $9$.
\end{warnbox}

% ---------------------------------------------------------------------
\shead{Dãy số có quy luật, trung bình cộng và trung vị}
\begin{knowbox}
\textbf{1. Dãy Fibonacci.} $1, 1, 2, 3, 5, 8, 13, 21, 34, 55, \ldots$ Mỗi số (kể từ số thứ ba) bằng \emph{tổng hai
số đứng ngay trước nó}: $2 = 1+1$, $3 = 1+2$, $5 = 2+3$, \ldots

\textbf{2. Cấp số cộng (dãy cách đều).} Mỗi số (kể từ số thứ hai) bằng số liền trước cộng với một số không đổi $d$
(công sai). Ví dụ: $1, 4, 7, 10$ (với $d = 3$). Số thứ $n$ bằng $a_1 + (n-1)\times d$.

\textbf{3. Trung bình cộng.} Trung bình cộng của $n$ số $=$ (tổng các số) $: n$. Suy ra
\[ \text{Tổng} = \text{trung bình cộng} \times \text{số các số}. \]
\textbf{4. Phương pháp ``độ lệch''.} Tổng các độ lệch so với trung bình cộng luôn bằng $0$: phần ``thừa'' của các số
lớn hơn trung bình đúng bằng phần ``thiếu'' của các số nhỏ hơn trung bình.

\textbf{5. Trung vị.} Sắp xếp các số theo thứ tự tăng dần; số đứng chính giữa là trung vị (nếu có một số lẻ các số).
Ví dụ: trung vị của $42, 43, 46, 49, 50$ là $46$. Nếu có một số chẵn các số thì trung vị là trung bình cộng của
hai số đứng giữa: trung vị của $6, 7, 9, 10$ là $(7+9):2 = 8$. Lưu ý: khi thêm hoặc thay đổi một số, vị trí ``ở giữa''
có thể thay đổi, nên phải sắp xếp lại.

\textbf{6. Các số nguyên liên tiếp.} Với một số lẻ các số nguyên liên tiếp, trung bình cộng chính là số đứng giữa:
trung bình của $19, 20, 21, 22, 23$ là $21$.
\end{knowbox}
\begin{proofbox}
Vì sao tổng các độ lệch bằng $0$? Giả sử $4$ số có trung bình $80$, tức tổng là $4\times 80 = 320$. Viết mỗi số thành
$80 + $ (độ lệch). Cộng lại: $320 + $ (tổng các độ lệch) $= 320$, nên tổng các độ lệch phải bằng $0$.
\end{proofbox}
\begin{figC}
  \draw[Slate!60] (-0.3,0)--(5.3,0);
  \node[left,font=\scriptsize,text=Slate] at (-0.3,0) {TB $=40$};
  \foreach \x/\v/\c in {0/38/tmOrange,1.5/43/tmGreen,3/41/tmGreen,4.5/38/tmOrange}{
    \pgfmathsetmacro\h{(\v-40)*0.3}
    \fill[\c!70] (\x,0) rectangle (\x+0.7,\h);
    \node[font=\scriptsize] at (\x+0.35,{\h+0.25*sign(\h+0.001)}) {$\v$};}
\end{figC}
\begin{center}\small Bốn số $38, 43, 41, 38$ có trung bình cộng $40$. Độ lệch so với $40$: $-2,\ +3,\ +1,\ -2$; tổng $= 0$.\end{center}
\begin{parentbox}
Có thể minh họa trung bình cộng bằng các chồng xu: ``san bằng'' các chồng cao xuống chồng thấp cho đến khi các chồng
cao bằng nhau. Chiều cao chung đó là trung bình cộng. Số xu chuyển đi từ chồng cao đúng bằng số xu thêm vào chồng thấp
-- đó chính là ý ``tổng độ lệch bằng $0$''.
\end{parentbox}

% =====================================================================
% VÍ DỤ MẪU
% =====================================================================
\needspace{14\baselineskip}
\par\vspace{0.7em}\noindent{\large\bfseries\color{tmOrange!90!black}\faPencil*\ \ Ví dụ mẫu}\par\vspace{0.3em}

\begin{exbox}{Ví dụ 1}
Cho tập hợp $\{0{,}5;\ 0{,}45;\ 0{,}405;\ 0{,}054;\ 0{,}54\}$.
\begin{enumerate}[label=\alph*),leftmargin=1.6em,itemsep=0pt,topsep=2pt]
  \item Sắp xếp các số theo thứ tự tăng dần.
  \item Tính tổng của số lớn nhất và số nhỏ nhất.
\end{enumerate}
\guide
\step{1}{Viết các số thành một cột, thêm chữ số $0$ cho đủ ba chữ số sau dấu phẩy}
Thêm chữ số $0$ vào tận cùng phần thập phân không làm thay đổi giá trị, nên ta có:
\[ 0{,}5 = 0{,}500;\quad 0{,}45 = 0{,}450;\quad 0{,}405;\quad 0{,}054;\quad 0{,}54 = 0{,}540. \]
\step{2}{So sánh như các số tự nhiên}
Bỏ ``$0{,}$'' đi, ta so sánh các số $500;\ 450;\ 405;\ 54;\ 540$:
\[ 54 < 405 < 450 < 500 < 540. \]
Vậy thứ tự tăng dần là: $0{,}054 < 0{,}405 < 0{,}45 < 0{,}5 < 0{,}54$.
\step{3}{Cộng số lớn nhất với số nhỏ nhất (đặt dấu phẩy thẳng cột)}
\[ \begin{array}{r} 0{,}540\\ +\ 0{,}054\\ \hline 0{,}594 \end{array} \]
\verify{$0{,}594 - 0{,}54 = 0{,}054$. Đúng.}
\ans{a) $0{,}054 < 0{,}405 < 0{,}45 < 0{,}5 < 0{,}54$;\quad b) $0{,}594$}
\end{exbox}

\begin{exbox}{Ví dụ 2}
Emma đạp xe trên quãng đường dài $45$ km. Buổi sáng bạn ấy đi được $\dfrac{2}{5}$ quãng đường. Hỏi buổi chiều Emma
còn phải đi bao nhiêu ki-lô-mét nữa, và quãng đường đó bằng mấy phần của cả quãng đường?
\guide
\step{1}{Tính quãng đường buổi sáng}
\begin{recallbox}Muốn tìm $\dfrac{m}{n}$ của số $a$, ta tính $a\times\dfrac{m}{n}$.\end{recallbox}
Buổi sáng: $45 \times \dfrac{2}{5} = \dfrac{45\times 2}{5} = \dfrac{90}{5} = 18$ (km).
\step{2}{Gọi ẩn và lập phương trình}
Gọi $x$ (km) là quãng đường buổi chiều còn phải đi ($x > 0$). Quãng đường buổi sáng cộng buổi chiều bằng cả quãng đường:
\[ 18 + x = 45. \]
\step{3}{Giải phương trình}
$x = 45 - 18 = 27$ (km).
\step{4}{Phần quãng đường còn lại}
Cách 1: $1 - \dfrac{2}{5} = \dfrac{5}{5} - \dfrac{2}{5} = \dfrac{3}{5}$.\quad
Cách 2: $\dfrac{27}{45} = \dfrac{27:9}{45:9} = \dfrac{3}{5}$. Hai cách cho cùng một kết quả.
\verify{$18 + 27 = 45$ km và $45\times\dfrac35 = 27$ km. Đúng.}
\ans{$27$ km, bằng $\dfrac{3}{5}$ quãng đường}
\end{exbox}

\begin{exbox}{Ví dụ 3}
a) Tính $20 - 4\times 3 + 8 : 2$ \quad và \quad $(20-4)\times(3+8) : 2$.\\
b) Không tính chính xác, hãy cho biết $101 \times 49$ gần nhất với số nào: $50$; $500$; $5000$; $50\,000$? Sau đó
tính chính xác bằng cách nhanh.
\guide
\step{1}{Biểu thức không có ngoặc: nhân, chia trước; cộng, trừ sau}
\[ 20 - \underbrace{4\times 3}_{12} + \underbrace{8 : 2}_{4} = 20 - 12 + 4 = 8 + 4 = 12. \]
(Phép trừ và phép cộng làm lần lượt từ trái sang phải: $20 - 12 = 8$, rồi $8 + 4 = 12$.)
\step{2}{Biểu thức có ngoặc: làm trong ngoặc trước}
\[ (20 - 4)\times(3+8) : 2 = 16 \times 11 : 2 = 176 : 2 = 88. \]
\step{3}{Ước lượng}
$101 \approx 100$ và $49 \approx 50$, nên $101 \times 49 \approx 100\times 50 = 5000$. Chọn $5000$.
\step{4}{Tính chính xác bằng tính chất phân phối}
$101 \times 49 = (100 + 1)\times 49 = 4900 + 49 = 4949$, quả thật rất gần $5000$.
\ans{a) $12$ và $88$;\quad b) gần $5000$ nhất; $101\times 49 = 4949$}
\end{exbox}

\begin{exbox}{Ví dụ 4}
a) Táo giá $\$1{,}80$ một ki-lô-gam. Grace mua $15$ kg và được giảm $10\%$ trên tổng hóa đơn. Grace phải trả
bao nhiêu tiền?\\
b) Một chiếc áo giá $\$40$. Cứ sau mỗi tuần, cửa hàng lại giảm $10\%$ \emph{giá hiện tại}. Sau ít nhất bao
nhiêu tuần thì giá áo thấp hơn $\$30$?
\guide
\step{1}{Thành tiền khi chưa giảm}
Thành tiền $=$ đơn giá $\times$ số lượng $= 1{,}80 \times 15 = 27$ (đô la).
\step{2}{Trừ đi $10\%$}
$10\%$ của $27$ là $27 : 10 = 2{,}70$ (đô la). Grace trả $27 - 2{,}70 = 24{,}30$ (đô la).
(Hoặc: còn $90\%$, tức $27 \times 0{,}9 = 24{,}3$.)
\step{3}{Giảm giá liên tiếp: mỗi tuần nhân với $0{,}9$}
Gọi $g_n$ là giá áo (đô la) sau $n$ tuần. Ta có $g_0 = 40$ và $g_{n} = g_{n-1}\times 0{,}9$.
\begin{center}
\begin{tabular}{|c|c|c|c|c|}
\hline
\rowcolor{tmBlue!10} Thời điểm & Ban đầu & Sau 1 tuần & Sau 2 tuần & Sau 3 tuần\\ \hline
Giá (đô la) & $40$ & $36$ & $32{,}40$ & $29{,}16$\\ \hline
Được giảm (đô la) & -- & $4$ & $3{,}60$ & $3{,}24$\\ \hline
\end{tabular}
\end{center}
Sau $2$ tuần giá là $\$32{,}40 > \$30$; sau $3$ tuần giá là $\$29{,}16 < \$30$.
\begin{warnbox}
Không được nghĩ ``mỗi tuần giảm $\$4$'' (vì $10\%$ của $40$ là $4$). Mỗi tuần giảm $10\%$ của \emph{giá lúc đó}, nên số
tiền được giảm ngày càng ít: $4$, rồi $3{,}60$, rồi $3{,}24$, \ldots
\end{warnbox}
\ans{a) $\$24{,}30$;\quad b) $3$ tuần}
\end{exbox}

\begin{exbox}{Ví dụ 5}
a) Leo đạp xe $3$ vòng quanh một sân vận động có chu vi $400$ m hết $4$ phút. Mia đạp xe với vận tốc $15$ km/h.
Ai đi nhanh hơn?\\
b) Một ô tô chạy với vận tốc $72$ km/h. Trong $3$ giây, ô tô đi được bao nhiêu mét?
\guide
\step{1}{Vận tốc của Leo (m/s)}
Quãng đường: $3 \times 400 = 1200$ (m). Thời gian: $4$ phút $= 4\times 60 = 240$ (giây).
\[ v_{\text{Leo}} = 1200 : 240 = 5 \text{ (m/s)}. \]
\step{2}{Đổi vận tốc của Mia sang m/s}
\begin{recallbox}$1$ km/h $= \dfrac{1000 \text{ m}}{3600 \text{ s}} = \dfrac{5}{18}$ m/s.\end{recallbox}
$v_{\text{Mia}} = 15 \times \dfrac{5}{18} = \dfrac{75}{18} \approx 4{,}17$ (m/s). Vì $5 > 4{,}17$ nên \textbf{Leo nhanh hơn}.
(Cách khác: $5$ m/s $= 5\times 3{,}6 = 18$ km/h $> 15$ km/h.)
\step{3}{Ô tô}
$72$ km/h $= 72 \times \dfrac{5}{18} = 20$ m/s. Trong $3$ giây đi được $20 \times 3 = 60$ (m).
\ans{a) Leo nhanh hơn ($5$ m/s so với khoảng $4{,}17$ m/s);\quad b) $60$ m}
\end{exbox}

\begin{exbox}{Ví dụ 6}
a) Trong các số $72;\ 84;\ 117;\ 135$, số nào chia hết cho $9$?\\
b) Tìm chữ số $x$ để số $\nd{3x5}$ chia hết cho $9$.\\
c) Có bao nhiêu số tự nhiên là ước của cả $24$ và $36$?
\guide
\step{1}{Dùng dấu hiệu chia hết cho $9$: tính tổng các chữ số}
$72$: $7+2 = 9$ (chia hết cho $9$). \quad $84$: $8 + 4 = 12$ (không chia hết cho $9$).\\
$117$: $1+1+7 = 9$ (chia hết). \quad $135$: $1+3+5 = 9$ (chia hết).\\
Vậy $72$, $117$, $135$ chia hết cho $9$; $84$ thì không ($84 = 9\times 9 + 3$).
\step{2}{Tìm chữ số $x$}
Tổng các chữ số của $\nd{3x5}$ là $3 + x + 5 = 8 + x$. Vì $x$ là chữ số nên $8 \le 8 + x \le 17$. Trong khoảng này chỉ có
$9$ chia hết cho $9$, nên $8 + x = 9$, tức $x = 1$. Số cần tìm là $315 = 9 \times 35$.
\step{3}{Ước chung của $24$ và $36$}
Liệt kê ước của $24$: $1, 2, 3, 4, 6, 8, 12, 24$. Ước của $36$: $1, 2, 3, 4, 6, 9, 12, 18, 36$.\\
Các số có mặt ở cả hai danh sách: $1, 2, 3, 4, 6, 12$ -- đó cũng chính là các ước của ƯCLN$(24,36) = 12$.
\ans{a) $72;\ 117;\ 135$;\quad b) $x = 1$;\quad c) $6$ số}
\end{exbox}

\begin{exbox}{Ví dụ 7}
Trong phép tính $\nd{AB} + \nd{AB} + \nd{AB} = \nd{1AB}$, mỗi chữ cái đại diện cho một chữ số. Tìm $A$ và $B$.
\guide
\textbf{Cách 1 -- Dùng cấu tạo số.} Ta có $\nd{1AB} = 100 + \nd{AB}$. Vậy
\[ 3\times\nd{AB} = 100 + \nd{AB} \;\Rightarrow\; 2\times\nd{AB} = 100 \;\Rightarrow\; \nd{AB} = 50. \]
Vậy $A = 5$, $B = 0$.

\textbf{Cách 2 -- Xét từng hàng.} Hàng đơn vị: $B + B + B = 3B$ có chữ số tận cùng là $B$, tức $3B - B = 2B$ tận cùng
bằng $0$, nên $B = 0$ hoặc $B = 5$.
\begin{itemize}[leftmargin=1.4em,itemsep=1pt,topsep=1pt]
  \item Nếu $B = 5$: $3\times 5 = 15$, viết $5$ nhớ $1$. Hàng chục: $3A + 1 = \nd{1A} = 10 + A$, suy ra $2A = 9$ -- không có
  chữ số nào thỏa mãn. Loại.
  \item Nếu $B = 0$: không có nhớ. Hàng chục: $3A = 10 + A$, suy ra $2A = 10$, $A = 5$.
\end{itemize}
\verify{$50 + 50 + 50 = 150 = \nd{1AB}$ với $A = 5$, $B = 0$. Đúng.}
\ans{$A = 5$, $B = 0$}
\end{exbox}

\begin{exbox}{Ví dụ 8}
a) Bốn bài kiểm tra của Noah được $6$; $9$; $10$ và $x$ điểm. Điểm trung bình của bốn bài là $8$. Tìm $x$.\\
b) Trung bình cộng của năm số tự nhiên liên tiếp là $30$. Tìm số nhỏ nhất.
\guide
\step{1}{Lập phương trình từ định nghĩa trung bình cộng}
\begin{recallbox}Trung bình cộng $=$ tổng $:$ số các số, nên tổng $=$ trung bình cộng $\times$ số các số.\end{recallbox}
\[ \frac{6 + 9 + 10 + x}{4} = 8 \;\Rightarrow\; 25 + x = 32 \;\Rightarrow\; x = 32 - 25 = 7. \]
\step{2}{Cách dùng độ lệch (nhanh hơn)}
Độ lệch so với $8$: $-2$, $+1$, $+2$, tổng $+1$. Tổng mọi độ lệch phải bằng $0$, nên $x = 8 - 1 = 7$.
\step{3}{Năm số liên tiếp}
Gọi số nhỏ nhất là $n$; năm số là $n, n+1, n+2, n+3, n+4$, có tổng $5n + 10$. Ta có
$5n + 10 = 5\times 30 = 150$, nên $5n = 140$, $n = 28$.
Năm số là $28, 29, 30, 31, 32$; số nhỏ nhất là $28$.
\verify{$28 + 29 + 30 + 31 + 32 = 150 = 5\times 30$. Đúng.}
\ans{a) $x = 7$;\quad b) $28$}
\end{exbox}

% =====================================================================
\clearpage
\dividerbar{Luyện tập cơ bản}

% ---------------------------------------------------------------- Bài 1
\needspace{10\baselineskip}
\begin{probbox}{1}{\relax}
Số nhỏ nhất trong tập hợp $\{3{,}2;\ 2{,}3;\ 3;\ 2{,}23;\ 3{,}22\}$ là
\choices{$3{,}2$}{$2{,}3$}{$3$}{$2{,}23$}{$3{,}22$}
\end{probbox}
\probsol{
\textbf{Phân tích đề.} Có $5$ số thập phân. Ta cần tìm số \emph{nhỏ nhất}.
\begin{recallbox}So sánh số thập phân: so sánh phần nguyên trước; nếu phần nguyên bằng nhau thì so sánh lần lượt hàng
phần mười, hàng phần trăm, \ldots\ Có thể thêm chữ số $0$ vào cuối phần thập phân mà giá trị không đổi.\end{recallbox}
\step{1}{So sánh phần nguyên}
Phần nguyên của các số: $3{,}2 \to 3$; \ $2{,}3 \to 2$; \ $3 \to 3$; \ $2{,}23 \to 2$; \ $3{,}22 \to 3$.\\
Số nhỏ nhất phải có phần nguyên nhỏ nhất, tức là bằng $2$. Chỉ còn hai ứng viên: $2{,}3$ và $2{,}23$.
\step{2}{Viết cùng số chữ số thập phân và so sánh}
$2{,}3 = 2{,}30$ \ và \ $2{,}23$. Hàng phần mười: $3$ so với $2$; vì $2 < 3$ nên $2{,}23 < 2{,}30$.
\step{3}{Kết luận}
Sắp xếp toàn bộ cho chắc chắn: $2{,}23 < 2{,}3 < 3 < 3{,}2 < 3{,}22$. Số nhỏ nhất là $2{,}23$.
\begin{warnbox}Đừng chọn $2{,}3$ chỉ vì ``$3$ nhỏ hơn $23$''. Chữ số $3$ trong $2{,}3$ là $3$ phần mười, còn trong
$2{,}23$ chỉ có $2$ phần mười.\end{warnbox}
\verify{$2{,}23 = \dfrac{223}{100}$ và $2{,}3 = \dfrac{230}{100}$; $223 < 230$.}
Số $2{,}23$ là phương án \textbf{(D)}.
\ans{(D) $2{,}23$}
}

% ---------------------------------------------------------------- Bài 2
\needspace{10\baselineskip}
\begin{probbox}{2}{\relax}
Keegan tham gia một chuyến chèo thuyền kayak dài $36$ km. Trước bữa trưa, Keegan đã chèo được $12$ km đầu tiên. Hỏi
sau bữa trưa, Keegan còn phải chèo bao nhiêu phần của cả chuyến đi?
\choices{$\dfrac{1}{2}$}{$\dfrac{5}{6}$}{$\dfrac{3}{4}$}{$\dfrac{2}{3}$}{$\dfrac{3}{5}$}
\end{probbox}
\probsol{
\textbf{Phân tích đề.} Cả chuyến dài $36$ km, đã chèo $12$ km. Đề hỏi phần \emph{còn lại} (chưa chèo), viết dưới dạng phân số.
\step{1}{Tính quãng đường còn lại}
Gọi $x$ (km) là quãng đường còn phải chèo. Ta có $12 + x = 36$, nên $x = 36 - 12 = 24$ (km).
\step{2}{Lập phân số ``phần còn lại trên toàn bộ''}
\begin{recallbox}Phân số chỉ ``bao nhiêu phần của toàn thể'' $=\dfrac{\text{phần đang xét}}{\text{toàn thể}}$. Rút gọn bằng cách
chia cả tử và mẫu cho ước chung lớn nhất.\end{recallbox}
\[ \frac{24}{36} = \frac{24 : 12}{36 : 12} = \frac{2}{3}. \]
\step{3}{Cách khác: tính phần đã chèo rồi lấy $1$ trừ đi}
Phần đã chèo: $\dfrac{12}{36} = \dfrac{1}{3}$. Phần còn lại: $1 - \dfrac13 = \dfrac33 - \dfrac13 = \dfrac23$. Hai cách cho cùng kết quả.
\begin{warnbox}Phân số $\dfrac{12}{36} = \dfrac{1}{3}$ là phần \emph{đã chèo}. Đề hỏi phần \emph{còn lại}, nên đáp án là
$1 - \dfrac13 = \dfrac23$.\end{warnbox}
\verify{$\dfrac{2}{3}\times 36 = 24$ km; $24 + 12 = 36$ km. Đúng.}
\begin{parentbox}Hãy hỏi con hai câu: ``Đã làm được bao nhiêu phần?'' và ``Còn lại bao nhiêu phần?''. Tập thói quen
trả lời cả hai câu giúp con tránh được bẫy chọn nhầm $\dfrac13$.\end{parentbox}
Phân số $\dfrac{2}{3}$ là phương án \textbf{(D)}.
\ans{(D) $\dfrac{2}{3}$}
}

% ---------------------------------------------------------------- Bài 3
\needspace{10\baselineskip}
\begin{probbox}{3}{\relax}
Khẳng định nào sau đây là \textbf{đúng}?
\choicesII{$0$ nhỏ hơn $-5$}{$7$ nhỏ hơn $-1$}{$10$ nhỏ hơn $\dfrac{1}{4}$}{$-1$ nhỏ hơn $-3$}{$-8$ nhỏ hơn $-2$}
\end{probbox}
\probsol{
\textbf{Phân tích đề.} Mỗi phương án là một phép so sánh. Ta kiểm tra từng phương án, tìm phương án đúng.
\begin{recallbox}Trên trục số nằm ngang, số bên trái nhỏ hơn số bên phải. Mọi số âm nhỏ hơn $0$ và nhỏ hơn mọi số dương.
Trong hai số âm, số nào xa $0$ hơn thì nhỏ hơn.\end{recallbox}
\step{1}{Đặt các số lên trục số}
\begin{figC}
  \draw[->,thick,Navy] (-9.4*0.6,0)--(10.6*0.6,0);
  \foreach \x in {-9,...,10}{\draw[Navy] (\x*0.6,0.08)--(\x*0.6,-0.08);}
  \foreach \x in {-8,-5,-3,-2,-1,0,7,10}{\node[below,font=\scriptsize] at (\x*0.6,-0.1) {$\x$};}
  \foreach \x in {-8,-2}{\fill[tmOrange] (\x*0.6,0) circle (3pt);}
  \draw[->,tmGreen!70!black,thick,shorten <=3.5pt,shorten >=3.5pt] (-8*0.6,0) to[bend left=50] node[above,font=\scriptsize]{sang phải $6$ đơn vị} (-2*0.6,0);
\end{figC}
\step{2}{Kiểm tra từng phương án}
(A) $0$ nhỏ hơn $-5$? Sai: $-5$ nằm bên trái $0$, nên $-5 < 0$.\\
(B) $7$ nhỏ hơn $-1$? Sai: số dương luôn lớn hơn số âm.\\
(C) $10$ nhỏ hơn $\dfrac14$? Sai: $\dfrac14 = 0{,}25 < 10$.\\
(D) $-1$ nhỏ hơn $-3$? Sai: $-3$ xa $0$ hơn nên $-3 < -1$.\\
(E) $-8$ nhỏ hơn $-2$? \textbf{Đúng}: $-8$ nằm bên trái $-2$ (xa $0$ hơn).
\begin{warnbox}Nhiều bạn chọn (D) vì nghĩ ``$1 < 3$ nên $-1 < -3$''. Với số âm thì ngược lại: $-1 > -3$
(nhiệt độ $-1^\circ$C ấm hơn $-3^\circ$C).\end{warnbox}
\verify{$-2 - (-8) = 6 > 0$, nghĩa là $-2$ lớn hơn $-8$.}
Khẳng định đúng duy nhất là \textbf{(E)}.
\ans{(E) $-8$ nhỏ hơn $-2$}
}

% ---------------------------------------------------------------- Bài 4
\needspace{10\baselineskip}
\begin{probbox}{4}{\relax}
Giá trị của biểu thức $6 \times 2 - 3$ là
\choices{$9$}{$-6$}{$12$}{$15$}{$10$}
\end{probbox}
\probsol{
\textbf{Phân tích đề.} Biểu thức có hai phép tính: nhân và trừ, không có ngoặc.
\begin{recallbox}Trong biểu thức không có ngoặc: nhân, chia trước; cộng, trừ sau.\end{recallbox}
\step{1}{Thực hiện phép nhân trước}
$6\times 2 = 12$.
\step{2}{Rồi thực hiện phép trừ}
$12 - 3 = 9$.
\begin{warnbox}Nếu làm phép trừ trước, ta sẽ được $6\times(2-3) = 6\times(-1) = -6$ -- đó là phương án (B), một ``bẫy''
dành cho bạn quên thứ tự thực hiện phép tính. Biểu thức đề bài \emph{không có ngoặc}, nên phải nhân trước.\end{warnbox}
\verify{Đọc biểu thức như một câu chuyện: ``$6$ hộp, mỗi hộp $2$ cái bánh, ăn mất $3$ cái'' thì còn $12 - 3 = 9$ cái.}
Giá trị $9$ là phương án \textbf{(A)}.
\ans{(A) $9$}
}

% ---------------------------------------------------------------- Bài 5
\needspace{10\baselineskip}
\begin{probbox}{5}{\relax}
Bài kiểm tra có tổng điểm là $50$. Mukesh đạt $80\%$ số điểm. Mukesh được bao nhiêu điểm?
\choices{$40$}{$62{,}5$}{$10$}{$45$}{$35$}
\end{probbox}
\probsol{
\textbf{Phân tích đề.} Tổng điểm tối đa là $50$; Mukesh đạt $80\%$ của $50$. Cần tìm số điểm của Mukesh.
\begin{recallbox}Muốn tìm $p\%$ của số $a$ ta tính $a \times \dfrac{p}{100}$.\end{recallbox}
\step{1}{Viết phép tính}
$80\%$ của $50$ là $50 \times \dfrac{80}{100}$.
\step{2}{Tính}
$50 \times \dfrac{80}{100} = \dfrac{50\times 80}{100} = \dfrac{4000}{100} = 40$.
\step{3}{Cách nhẩm nhanh}
$10\%$ của $50$ là $5$; vậy $80\% = 8\times 10\%$ là $8\times 5 = 40$.
\begin{warnbox}Phương án (B) $62{,}5$ là kết quả của phép chia $50 : 0{,}8$ -- sai, vì $62{,}5$ còn lớn hơn cả điểm tối đa $50$!
Phương án (C) $10$ là số điểm Mukesh \emph{bị mất} ($20\%$ của $50$).\end{warnbox}
\verify{$\dfrac{40}{50} = \dfrac{80}{100} = 80\%$. Đúng.}
Số điểm $40$ là phương án \textbf{(A)}.
\ans{(A) $40$}
}

% ---------------------------------------------------------------- Bài 6
\needspace{10\baselineskip}
\begin{probbox}{6}{\relax}
Thịt bò xay được bán với giá $\$5{,}00$ một ki-lô-gam. Mua $12$ kg thịt bò xay thì hết bao nhiêu tiền?
\choices{$\$5{,}00$}{$\$12{,}00$}{$\$60{,}00$}{$\$17{,}00$}{$\$2{,}40$}
\end{probbox}
\probsol{
\textbf{Phân tích đề.} Đơn giá (giá của $1$ kg) là $\$5{,}00$; số lượng là $12$ kg. Cần tìm thành tiền.
\begin{recallbox}Thành tiền $=$ đơn giá $\times$ số lượng.\end{recallbox}
\step{1}{Viết phép tính}
Gọi $T$ (đô la) là số tiền phải trả. Khi đó $T = 5 \times 12$.
\step{2}{Tính}
$T = 5\times 12 = 60$ (đô la), tức là $\$60{,}00$.
\step{3}{Ước lượng để kiểm tra}
Mỗi ki-lô-gam giá $\$5$, mua hơn $10$ kg nên phải trả hơn $5\times 10 = 50$ đô la; chỉ (C) thỏa mãn.
\begin{warnbox}Các phương án sai đến từ việc dùng nhầm phép tính: (A) là giá của $1$ kg; (B) chỉ là số ki-lô-gam viết thành
tiền; (D) là $5 + 12$ (cộng thay vì nhân); (E) là $12 : 5$ (chia thay vì nhân).\end{warnbox}
\verify{$60 : 12 = 5$ đô la mỗi ki-lô-gam. Đúng.}
Số tiền $\$60{,}00$ là phương án \textbf{(C)}.
\ans{(C) $\$60{,}00$}
}

% ---------------------------------------------------------------- Bài 7
\needspace{10\baselineskip}
\begin{probbox}{7}{\relax}
Giá trị của biểu thức $(8 \times 4) + 3$ là
\choices{$96$}{$15$}{$56$}{$35$}{$28$}
\end{probbox}
\probsol{
\textbf{Phân tích đề.} Biểu thức có ngoặc chứa phép nhân, sau đó cộng thêm $3$.
\begin{recallbox}Biểu thức có ngoặc: làm phép tính trong ngoặc trước.\end{recallbox}
\step{1}{Tính trong ngoặc}
$8\times 4 = 32$.
\step{2}{Cộng thêm $3$}
$32 + 3 = 35$.
\begin{warnbox}Đọc nhầm thành $8 \times (4 + 3) = 8\times 7 = 56$ sẽ dẫn đến phương án (C). Ngoặc trong đề bao quanh
$8\times 4$, không phải $4 + 3$. Phương án (A) $96 = 8\times 4\times 3$ là do nhầm dấu cộng thành dấu nhân.\end{warnbox}
\verify{Không có ngoặc thì kết quả vẫn là $8\times 4 + 3 = 35$, vì nhân luôn được làm trước cộng; ngoặc ở đây chỉ để nhấn mạnh.}
Giá trị $35$ là phương án \textbf{(D)}.
\ans{(D) $35$}
}

% ---------------------------------------------------------------- Bài 8
\needspace{10\baselineskip}
\begin{probbox}{8}{\relax}
Số nguyên nào sau đây gần với $99 \times 9$ nhất?
\choices{$10\,000$}{$100$}{$100\,000$}{$1000$}{$10$}
\end{probbox}
\probsol{
\textbf{Phân tích đề.} Các phương án là $10$; $100$; $1000$; $10\,000$; $100\,000$ -- cách nhau gấp $10$ lần. Chỉ cần ước
lượng $99\times 9$ là đủ, nhưng ta cũng tính chính xác để chắc chắn.
\step{1}{Ước lượng}
$99 \approx 100$ nên $99 \times 9 \approx 100\times 9 = 900$.
\step{2}{Tính chính xác bằng tính chất phân phối}
\begin{recallbox}Tính chất phân phối: $(a - b)\times c = a\times c - b\times c$.\end{recallbox}
$99\times 9 = (100 - 1)\times 9 = 900 - 9 = 891$.
\step{3}{So sánh khoảng cách đến các phương án gần nhất}
Khoảng cách từ $891$ đến $1000$ là $1000 - 891 = 109$.\\
Khoảng cách từ $891$ đến $100$ là $891 - 100 = 791$.\\
Các phương án còn lại ($10$; $10\,000$; $100\,000$) còn xa hơn nhiều. Vậy $1000$ là gần nhất.
\begin{figC}
  \draw[->,thick,Navy] (-0.3,0)--(10.6,0);
  \foreach \x/\l in {0/0,1/100,10/1000}{\draw[Navy] (\x,0.1)--(\x,-0.1); \node[below,font=\scriptsize] at (\x,-0.1) {$\l$};}
  \fill[tmOrange] (8.91,0) circle (3pt); \node[above,font=\scriptsize] at (8.91,0.1) {$891$};
  \draw[Slate!60,densely dotted] (1,-0.5)--(1,-0.95) (8.91,-0.1)--(8.91,-0.95) (10,-0.5)--(10,-0.95);
  \draw[<->,tmGreen!70!black] (8.91,-0.85)--(10,-0.85) node[midway,below,font=\scriptsize]{$109$};
  \draw[<->,AnsRed!70] (1,-0.85)--(8.91,-0.85) node[midway,below,font=\scriptsize]{$791$};
\end{figC}
\verify{$891 : 9 = 99$. Đúng.}
Số $1000$ là phương án \textbf{(D)}.
\ans{(D) $1000$}
}

% ---------------------------------------------------------------- Bài 9
\needspace{10\baselineskip}
\begin{probbox}{9}{\relax}
Joshua đang đọc một cuốn sách dày $396$ trang. Joshua mới chỉ đọc xong $\dfrac{1}{3}$ đầu tiên của cuốn sách. Hỏi
Joshua còn phải đọc bao nhiêu trang nữa để hết cuốn sách?
\choices{$264$}{$124$}{$250$}{$199$}{$244$}
\end{probbox}
\probsol{
\textbf{Phân tích đề.} Cả cuốn sách: $396$ trang. Đã đọc: $\dfrac13$ cuốn. Hỏi: số trang \emph{còn lại}.
\begin{recallbox}Muốn tìm $\dfrac{m}{n}$ của số $a$, ta tính $a\times\dfrac{m}{n}$ (chia $a$ cho $n$ rồi nhân với $m$).\end{recallbox}
\step{1}{Tính số trang đã đọc}
$\dfrac13$ của $396$ là $396\times\dfrac13 = 396 : 3 = 132$ (trang).
\step{2}{Gọi ẩn và lập phương trình}
Gọi $x$ là số trang Joshua còn phải đọc. Số trang đã đọc cộng số trang còn lại bằng cả cuốn sách:
\[ 132 + x = 396 \;\Rightarrow\; x = 396 - 132 = 264. \]
\step{3}{Cách khác: tính trực tiếp phần còn lại}
Phần còn lại là $1 - \dfrac13 = \dfrac23$ cuốn sách, tức $396 \times \dfrac23 = \dfrac{396\times 2}{3} = \dfrac{792}{3} = 264$ (trang).
\begin{warnbox}Dừng lại ở $132$ là trả lời nhầm câu hỏi (đó là số trang \emph{đã đọc}; và $132$ cũng không có trong các
phương án). Luôn đọc lại câu hỏi trước khi chọn đáp án.\end{warnbox}
\verify{$264 + 132 = 396$, và $132\times 3 = 396$. Đúng.}
Số $264$ là phương án \textbf{(A)}.
\ans{(A) $264$}
}

% ---------------------------------------------------------------- Bài 10
\needspace{10\baselineskip}
\begin{probbox}{10}{\relax}
Số nào sau đây là bội của $9$?
\choices{$50$}{$40$}{$35$}{$45$}{$55$}
\end{probbox}
\probsol{
\textbf{Phân tích đề.} Cần tìm số chia hết cho $9$ (là bội của $9$) trong năm số đã cho.
\begin{recallbox}Một số chia hết cho $9$ khi và chỉ khi tổng các chữ số của nó chia hết cho $9$.\end{recallbox}
\step{1}{Tính tổng các chữ số của từng số}
\begin{center}
\begin{tabular}{|c|c|c|c|c|c|}
\hline
\rowcolor{tmBlue!10} Số & $50$ & $40$ & $35$ & $45$ & $55$\\ \hline
Tổng các chữ số & $5$ & $4$ & $8$ & $9$ & $10$\\ \hline
Chia hết cho $9$? & không & không & không & \textbf{có} & không\\ \hline
\end{tabular}
\end{center}
\step{2}{Kiểm tra bằng bảng nhân $9$}
Các bội của $9$ gần đó: $36,\ 45,\ 54$. Chỉ có $45 = 9\times 5$ xuất hiện trong các phương án.
\begin{proofbox}Vì sao dấu hiệu đúng với số có hai chữ số $\nd{ab}$? Ta có $\nd{ab} = 10a + b = 9a + (a + b)$. Phần $9a$
chia hết cho $9$, nên $\nd{ab}$ chia hết cho $9$ đúng khi $a + b$ chia hết cho $9$. Với $45$: $45 = 9\times 4 + 9$.\end{proofbox}
\verify{$45 : 9 = 5$, không dư.}
Số $45$ là phương án \textbf{(D)}.
\ans{(D) $45$}
}

% ---------------------------------------------------------------- Bài 11
\needspace{10\baselineskip}
\begin{probbox}{11}{\relax}
Cho tập hợp $\{0{,}34;\ 0{,}304;\ 0{,}034;\ 0{,}43\}$. Tổng của số nhỏ nhất và số lớn nhất trong tập hợp là
\choices{$0{,}77$}{$0{,}734$}{$0{,}077$}{$0{,}464$}{$0{,}338$}
\end{probbox}
\probsol{
\textbf{Phân tích đề.} Bước 1: tìm số nhỏ nhất và số lớn nhất. Bước 2: cộng hai số đó.
\begin{recallbox}Để so sánh số thập phân, thêm chữ số $0$ cho các phần thập phân dài bằng nhau rồi so sánh từng hàng
từ trái sang phải.\end{recallbox}
\step{1}{Viết các số với ba chữ số thập phân}
\begin{center}
\begin{tabular}{|c|c|c|c|}
\hline
\rowcolor{tmBlue!10} Số ban đầu & Viết lại & Phần mười & Phần trăm \\ \hline
$0{,}34$ & $0{,}340$ & $3$ & $4$ \\ \hline
$0{,}304$ & $0{,}304$ & $3$ & $0$ \\ \hline
$0{,}034$ & $0{,}034$ & $0$ & $3$ \\ \hline
$0{,}43$ & $0{,}430$ & $4$ & $3$ \\ \hline
\end{tabular}
\end{center}
\step{2}{Tìm số nhỏ nhất và lớn nhất}
So sánh như số tự nhiên: $34 < 304 < 340 < 430$, nên $0{,}034 < 0{,}304 < 0{,}34 < 0{,}43$.\\
Số nhỏ nhất: $0{,}034$ (hàng phần mười là $0$). Số lớn nhất: $0{,}43$ (hàng phần mười là $4$).
\step{3}{Cộng, đặt dấu phẩy thẳng cột}
\[ \begin{array}{r} 0{,}430\\ +\ 0{,}034\\ \hline 0{,}464 \end{array} \]
\begin{warnbox}Đặt tính lệch hàng (cộng $43$ với $34$ ``theo hàng cuối'') sẽ cho $0{,}77$ hoặc $0{,}077$ -- phương án (A)
và (C). Phải để dấu phẩy thẳng cột: $4$ phần mười cộng $0$ phần mười, $3$ phần trăm cộng $3$ phần trăm, \ldots\end{warnbox}
\verify{$0{,}464 - 0{,}43 = 0{,}034$. Đúng.}
Tổng $0{,}464$ là phương án \textbf{(D)}.
\ans{(D) $0{,}464$}
}

% ---------------------------------------------------------------- Bài 12
\needspace{11\baselineskip}
\begin{probbox}{12}{\relax}
Keiko và Leah chạy trên một đường chạy vòng tròn dài $150$ m. Keiko chạy $3$ vòng hết $120$ giây, còn Leah chạy $5$ vòng hết
$160$ giây. Bạn nào chạy nhanh hơn, và vận tốc của bạn đó xấp xỉ bao nhiêu?
\choicesII{Keiko, $3{,}75$ m/s}{Keiko, $2{,}4$ m/s}{Leah, $3{,}3$ m/s}{Leah, $4{,}69$ m/s}{Leah, $3{,}75$ m/s}
\end{probbox}
\probsol{
\textbf{Phân tích đề.} Một vòng dài $150$ m. Keiko: $3$ vòng trong $120$ s. Leah: $5$ vòng trong $160$ s. Cần tính vận tốc
(m/s) của mỗi bạn, so sánh, rồi chọn phương án.
\begin{recallbox}Vận tốc $=$ quãng đường $:$ thời gian. Quãng đường chạy $n$ vòng $= n\times$ độ dài một vòng.\end{recallbox}
\step{1}{Quãng đường mỗi bạn chạy}
Keiko: $3\times 150 = 450$ (m). \quad Leah: $5\times 150 = 750$ (m).
\step{2}{Vận tốc của Keiko}
$v_{\text{Keiko}} = 450 : 120 = 3{,}75$ (m/s).
\step{3}{Vận tốc của Leah}
$v_{\text{Leah}} = 750 : 160 = 4{,}6875 \approx 4{,}69$ (m/s).
\step{4}{So sánh}
$4{,}69 > 3{,}75$ nên \textbf{Leah chạy nhanh hơn}, với vận tốc khoảng $4{,}69$ m/s.
\begin{center}
\begin{tabular}{|c|c|c|c|}
\hline
\rowcolor{tmBlue!10} & Quãng đường & Thời gian & Vận tốc\\ \hline
Keiko & $450$ m & $120$ s & $3{,}75$ m/s\\ \hline
Leah & $750$ m & $160$ s & $\approx 4{,}69$ m/s\\ \hline
\end{tabular}
\end{center}
\begin{warnbox}Phương án (E) ``Leah, $3{,}75$ m/s'' ghép đúng tên người nhanh hơn với vận tốc của người chậm hơn. Phương án (B)
$2{,}4$ m/s là kết quả của phép tính ngược $120 : 50$ (chia thời gian cho quãng đường). Hãy luôn viết rõ đơn vị:
m $:$ s $=$ m/s.\end{warnbox}
\verify{$4{,}6875 \times 160 = 750$ m $= 5$ vòng; $3{,}75\times 120 = 450$ m $= 3$ vòng. Đúng.}
Kết luận ``Leah, $4{,}69$ m/s'' là phương án \textbf{(D)}.
\ans{(D) Leah, $4{,}69$ m/s}
}

% ---------------------------------------------------------------- Bài 13
\needspace{16\baselineskip}
\begin{probbox}{13}{\relax}
Trong phép cộng dưới đây, $P$ và $Q$ là hai chữ số. Giá trị của $P + Q$ là
\[ \begin{array}{r@{\;}c@{\;}c@{\;}c} & P & Q & Q\\ & P & P & Q\\ + & Q & Q & Q\\ \hline & 8 & 7 & 6 \end{array} \]
\choices{$3$}{$5$}{$7$}{$6$}{$4$}
\end{probbox}
\probsol{
\textbf{Phân tích đề.} Ba số $\nd{PQQ}$, $\nd{PPQ}$, $\nd{QQQ}$ có tổng là $876$. Cần tìm hai chữ số $P$, $Q$ rồi tính $P + Q$.
Ta xét từng hàng, bắt đầu từ hàng đơn vị.
\step{1}{Hàng đơn vị: $Q + Q + Q = 3Q$ có chữ số tận cùng là $6$}
$Q$ là chữ số nên $3Q$ có thể là $0, 3, 6, \ldots, 27$. Các số tận cùng bằng $6$ trong đó: chỉ có $6$ (vì $16$ và $26$
không chia hết cho $3$). Vậy $3Q = 6$, $Q = 2$, và \emph{không có nhớ} sang hàng chục.
\step{2}{Hàng chục: $Q + P + Q = P + 4$ có chữ số tận cùng là $7$}
$P + 4$ nằm trong khoảng $4$ đến $13$; tận cùng bằng $7$ thì $P + 4 = 7$, tức $P = 3$. Không có nhớ sang hàng trăm.
\step{3}{Hàng trăm: kiểm tra $P + P + Q = 8$}
$3 + 3 + 2 = 8$. Khớp với chữ số hàng trăm của $876$.
\begin{recallbox}Cấu tạo số: $\nd{PQQ} = 100P + 11Q$, \ $\nd{PPQ} = 110P + Q$, \ $\nd{QQQ} = 111Q$.\end{recallbox}
\step{4}{Cách khác bằng cấu tạo số}
Cộng lại: $(100P + 11Q) + (110P + Q) + 111Q = 210P + 123Q = 876$. Thử $Q = 2$: $210P = 876 - 246 = 630$, nên $P = 3$.
(Các giá trị khác của $Q$ không cho $P$ là chữ số.)
\verify{$322 + 332 + 222 = 876$. Đúng.}
\begin{parentbox}Bài đố chữ giống một trò chơi trinh thám: mỗi hàng là một ``manh mối''. Hãy cho con viết riêng từng
hàng thành một câu (``ba lần $Q$ tận cùng bằng $6$'') rồi thử các chữ số từ $0$ đến $9$. Đừng quên kiểm tra số nhớ.\end{parentbox}
Vậy $P + Q = 3 + 2 = 5$, là phương án \textbf{(B)}.
\ans{(B) $5$}
}

% ---------------------------------------------------------------- Bài 14
\needspace{12\baselineskip}
\begin{probbox}{14}{\relax}
Trung bình cộng của bốn số nguyên $78$; $83$; $82$ và $x$ là $80$. Khẳng định nào sau đây đúng?
\choicesL{$x$ lớn hơn số trung bình cộng $2$ đơn vị}{$x$ nhỏ hơn số trung bình cộng $1$ đơn vị}{$x$ nhỏ hơn số trung bình cộng $2$ đơn vị}{$x$ nhỏ hơn số trung bình cộng $3$ đơn vị}{$x$ bằng số trung bình cộng}
\end{probbox}
\probsol{
\textbf{Phân tích đề.} Biết trung bình cộng của bốn số là $80$ và ba trong bốn số. Tìm $x$ rồi so sánh với $80$.
\begin{recallbox}Tổng các số $=$ trung bình cộng $\times$ số các số.\end{recallbox}
\step{1}{Lập phương trình}
Theo định nghĩa trung bình cộng: $\dfrac{78 + 83 + 82 + x}{4} = 80$, tức là $78 + 83 + 82 + x = 80\times 4 = 320$.
\step{2}{Giải phương trình tìm $x$}
$78 + 83 + 82 = 243$, nên $243 + x = 320$, suy ra $x = 320 - 243 = 77$.
\step{3}{So sánh với trung bình cộng}
$80 - 77 = 3$: $x$ nhỏ hơn trung bình cộng $3$ đơn vị.
\step{4}{Cách nhanh -- phương pháp độ lệch}
So với $80$: $78$ thiếu $2$; $83$ thừa $3$; $82$ thừa $2$. Tổng độ lệch của ba số: $-2 + 3 + 2 = +3$.
Vì tổng tất cả các độ lệch phải bằng $0$, nên $x$ phải ``thiếu'' $3$: $x = 80 - 3 = 77$.
\begin{proofbox}Mỗi số viết được thành $80 + $(độ lệch). Tổng bốn số là $4\times 80 + $(tổng độ lệch). Mà tổng bốn số đúng bằng
$4\times 80$, nên tổng độ lệch bằng $0$.\end{proofbox}
\verify{$(78 + 83 + 82 + 77) : 4 = 320 : 4 = 80$. Đúng.}
\begin{parentbox}Cách ``độ lệch'' rất hợp để tính nhẩm điểm trung bình trong đời sống. Có thể đố con: ``Ba bài
được $7$, $9$, $8$; muốn trung bình $8$ thì bài thứ tư cần mấy điểm?'' (Đáp: $8$, vì $-1 + 1 + 0 = 0$.)\end{parentbox}
Khẳng định ``$x$ nhỏ hơn số trung bình cộng $3$ đơn vị'' là phương án \textbf{(D)}.
\ans{(D)}
}

% ---------------------------------------------------------------- Bài 15
\needspace{10\baselineskip}
\begin{probbox}{15}{\relax}
Trong dãy Fibonacci $1, 1, 2, 3, 5, \ldots$, mỗi số kể từ số $2$ trở đi bằng tổng của hai số đứng ngay trước nó. Chẳng hạn,
số tiếp theo của dãy là $3 + 5 = 8$. Số nào sau đây có mặt trong dãy?
\choices{$20$}{$21$}{$22$}{$23$}{$24$}
\end{probbox}
\probsol{
\textbf{Phân tích đề.} Quy luật: số sau $=$ tổng hai số liền trước. Ta viết tiếp dãy cho đến khi vượt quá $24$ (số lớn
nhất trong các phương án), rồi xem phương án nào xuất hiện.
\step{1}{Viết tiếp dãy}
\begin{center}
\begin{tabular}{|c|c|c|c|c|c|c|c|c|c|}
\hline
\rowcolor{tmBlue!10} Thứ tự & 1 & 2 & 3 & 4 & 5 & 6 & 7 & 8 & 9\\ \hline
Số & $1$ & $1$ & $2$ & $3$ & $5$ & $8$ & $13$ & $\mathbf{21}$ & $34$\\ \hline
\end{tabular}
\end{center}
Cụ thể: $5 + 8 = 13$; \ $8 + 13 = 21$; \ $13 + 21 = 34$.
\step{2}{Đối chiếu với các phương án}
Số $34$ đã lớn hơn $24$, nên các số tiếp theo đều lớn hơn $24$. Giữa $13$ và $34$ chỉ có đúng một số của dãy là $21$.
Vậy $20$, $22$, $23$, $24$ không thuộc dãy; chỉ có $21$ thuộc dãy.
\begin{tipbox}Dãy Fibonacci luôn tăng (từ số thứ hai trở đi) vì ta cộng thêm số dương. Nhờ vậy, khi đã vượt qua một số,
ta biết chắc số đó không thể xuất hiện về sau nữa.\end{tipbox}
\verify{$8 + 13 = 21$. Đúng.}
Số $21$ là phương án \textbf{(B)}.
\ans{(B) $21$}
}
